\documentclass{article}
\usepackage{tikz}
\usetikzlibrary{angles, quotes}
\usepackage{amssymb}
\usepackage{pgfplots}
\pgfplotsset{width=10cm,compat=1.9}
\usepgfplotslibrary{statistics}
\usepackage[left=2.54cm,right=2.54cm,top=2.54cm,bottom=2.54cm]{geometry}
\usepackage{tasks} %Paquete necesario para la producción de items horizontales para algunos ejercicios de tarea empleando el comando /begin{multicols}}{#}.../end{multicols} para ayudarnos a reducir las páginas
\usepackage{pdflscape} %comando que permite cambiar a orientación horizontal las hojas%
\usepackage{soul} %Paquete para permitir la compilación de acentos usando el comando {\hl{}}}\\
\sethlcolor{yellow(munsell)} %comando para definir el color para el subrayado usando el comando \hl
\usepackage[utf8]{inputenc}
\usepackage[spanish]{babel}
\usepackage{icomma}
\usepackage{ragged2e}
\usepackage{siunitx}
\usepackage{url}
\usepackage[colorlinks=true, urlcolor=blue,  linkcolor=black, citecolor=green]{hyperref}
\usepackage{pdfpages}
\usepackage{blindtext} %Paquete que genera texto ficticio (dummy text) con el comando: \blindtext
\usepackage{cancel} %in the preamble gives you four different modes of striking through: \cancel{text to cancel} draws a diagonal line (slash) through its argument, \bcancel{text to cancel} uses the negative slope (a backslash), \xcancel{text to cancel} draws an X (actually \cancel plus \bcancel), \cancelto{〈value〉}{〈expression〉} draws a diagonal arrow through the 〈expression〉pointing to the 〈value〉 (math-mode only)
\usepackage{csquotes}
\usepackage{afterpage}
\usepackage{parskip} 
\usepackage{float}
\usepackage{enumitem}
\usepackage{multicol}%Paquete que permite la creación aislada de columnas de texto. De esta forma se reduce la cantidad de páginas en nuestro documento
\usepackage{lipsum} %Paquete que genera texto de relleno al igual que \usepackage{blindtext}, se usa con el comando \lipsum[#-#] los #(hashtags) delimitan el número de páginas que deseamos ocupar del paquete lipsum para usarlos como relleno. 

% Definir el comando personalizado para la nota
\newcommand{\mynote}[1]{%
\begin{tikzpicture}[baseline=-0.75ex]
\node[draw, circle, fill=Apple Green!60, inner sep=2pt] (note) {\textbf{Nota:}};
\end{tikzpicture}%
\ \textit{#1}%
}

\newenvironment{Figura}
{\par\medskip\noindent\minipage{\linewidth}}
{\endminipage\par\medskip}
\usepackage{caption}
\usepackage[
backend=biber,
sorting=none,
url=true
]{biblatex} %bibliografía
\addbibresource{biblio.bib}
% Establecer el espacio entre las entradas de la bibliografía
\setlength{\bibitemsep}{1\baselineskip} % Puedes ajustar el valor según tus preferencias
\usepackage{amsmath,amsthm,amssymb,amsfonts}
\usepackage{pifont} %Permite la compilación del comando \xmark: es el símbolo de la tachita
\usepackage{mathtools} %permite la compilación de símbolos para matrices
\usepackage{empheq} %Paquete que se relaciona con \usepackage[most]{tcolorbox} para la creación de cuadros/cajas de colores para delimitar los resultados matemáticos o líneas de texto usando la línea de comando: \begin{empheq}[box={\mymath[colback=orange(webcolor)!70,drop lifted shadow]}]{equation*} \end{empheq}
\usepackage[most]{tcolorbox} %Paquete que se relaciona con \usepackage{empheq} para la creación de cuadros/cajas de colores para delimitar los resultados matemáticos o líneas de texto
\renewcommand{\qedsymbol}{$\blacksquare$}
\usetikzlibrary{positioning,decorations.pathreplacing} %Paquete que permite la compilación de llaves para cuadros sinópticos (brace diagram) 
\usepackage{schemata} %The schemata package is designed just for these brace diagrams
\usepackage{booktabs}

\newcommand\AB[2]{\schema{\schemabox{#1}}{\schemabox{#2}}} %comando o paquete necesario para crear cuadros sinópticos

\newtcbox{\mymath}[1][]{%
nobeforeafter, math upper, tcbox raise base,
enhanced, colframe=blue!30!black,
colback=blue!30, boxrule=1pt,
#1} %Comando importante para encasillar los resultados matemáticos en cajas de diferentes colores, se relaciona con los paquetes \usepackage{empheq} y \usepackage[most]{tcolorbox} 

\newenvironment{sysmatrix}[1]
{\left(\begin{array}{@{}#1@{}}}
{\end{array}\right)}
\newcommand{\ro}[1]{%
\xrightarrow{\mathmakebox[\rowidth]{#1}}%
}
\newlength{\rowidth}% row operation width
\AtBeginDocument{\setlength{\rowidth}{3em}}  %Comando importante para laproducción de lineas de operaciones entre matrices, método de Gauss o Gauss Jordan se relaciona con \newenvironment{sysmatrix}

%formato para cambiar el horario a español
\usepackage[spanish]{datetime2}
\DTMsetdatestyle{spanish}

\renewcommand{\today}{\DTMdisplaydate{\the\year}{\the\month}{\the\day }{-1}}
%%%%%%%%%%%%%%%%%%%%%%%%%%%%%%%%%%%%%%%%%

\usepackage{datetime}
\newcommand{\mycurrenttime}{\xxivtime}

%%%%%%%%%%%%%%%%%%%%%%%%%%%%%%%%%%%%%%%%%%%%%%%%%%%%%%%%%%%%%%%%%%%%%%%%%%%%%%%%%%%%%%%%%%%%%%%%%%%%%%%%%%%%%%%%%%%%%%%%%%%%%%%%%%%%%%%%%%%%%%%%%


%%%%%%%%%%%%%Caja de color para el título%%%%%%%%%%%%%%%%%%%%%%%%%%%%%%%%%%%%%%%%%%%%%%%%%%%%%%

\definecolor{myframecolor}{RGB}{1, 45, 75} % Prussian Blue
\definecolor{myboxcolor}{RGB}{0, 107, 92} % Apple Green

\newtcolorbox{mybox}{
enhanced,
colback=myboxcolor!7, % Color de fondo del cuadro
colframe=myframecolor, % Color del marco del cuadro
arc=0pt,
boxrule=1pt,
borderline west={2mm}{-10mm}{myframecolor}, % Borde en el lado izquierdo
sharp corners=southwest,
width=\linewidth,
before=\par\vspace{\bigskipamount}, % Espacio antes del cuadro
after=\par\vspace{\bigskipamount} % Espacio después del cuadro
}


\newcommand{\euler}{e} %Comando para producir letra e de euler
\newenvironment{remark}{\par\vfill\footnotesize % Vertical white space above the remark and smaller font size
\begin{list}{}{
	\leftmargin=80pt % Indentation on the left
	\rightmargin=60pt}\item\ignorespaces % Indentation on the right
\makebox[-2.5pt]{\begin{tikzpicture}[overlay]
		\node[draw=Horizon!60,line width=2.5pt,circle,fill=Horizon!25,font=\sffamily\bfseries,inner sep=4pt,outer sep=0pt] at (-19pt,5pt){\textcolor{Horizon}{Nota}};\end{tikzpicture}} % Blue Nota in a circle
\advance\baselineskip -1pt}{\end{list}\vskip5pt} % Tighter line spacing and white space after remark
\usepackage{graphicx}

\usepackage{titling}

\input{colores.tex}

%%Producir signos de cita textual``Asignatura''

\newcommand{\xmark}{\ding{55}} %%%comando que genera la tachita

\newenvironment{MyColorPar}[1]{%
\leavevmode\color{#1}\ignorespaces%
}{%
}%

%%%%%%%%%%%%%%%%%%% Título de la tarea, Nombre de alumno

\title{Clase \# 36 - Mecánica Teórica}
\author{\emph{JA} $\boldsymbol{\mid}$}
\date{\today}

%%%%%%%%%%%%%%%%%%% Datos de la Materia

\usepackage{fancyhdr}
\fancypagestyle{plain}{%  the preset of fancyhdr 
\fancyhf{} % clear all header and footer fields
\fancyfoot[R]{
\includegraphics[width=2cm]{LogoFCFMBUAP (1).png}}
\fancyfoot[L]{{\bfseries{\thedate{}}} a las {\bfseries{\mycurrenttime{}}} horas{} (GMT-6, H. Puebla de Zaragoza, Pue)}
\fancyhead[L]{Mecánica teórica}
\fancyhead[R]{\theauthor}
}
\makeatletter
\def\@maketitle{%
\newpage
\null
\vskip 1em%
\begin{center}%
\let \footnote \thanks
{\LARGE \@title \par}%
\vskip 1em%
%{\large \@date}%
\end{center}%
\par
\vskip 1em}
\makeatother

\usepackage{lipsum}  
\usepackage{cmbright}


%%%%%%%%%%%%%%%%%%%%%%%%%%%%%%%%%%%%%%%%%%%%%%%%%%%%%%%%%%%%%%%%%%%%%%%%%%%%%%%%%%%%%%%%%%%%%%%%%%%%%%%%%%%%%%%%%%%%%%%%%%%%%%%%%%%%%%%%%%%%%%%%%%%%%%%%%%%%%%%%%%%%%%%%%%%%%%%%%%%%%%%%%%%%%%%%%%%%%%

\begin{document}



\maketitle


%%%%%%%%%%%%%%%%%%% Datos del profesor, campus y aula

\begin{mybox}
\noindent\begin{tabular}{@{}ll}
	{\bfseries{Profesor}} & Dr. Gilberto Silva Ortigoza\\
	\textcolor{prussianblue}{\bfseries{Campus}} / \textcolor{trueblue}{\bfseries{Facultad}}  & \textcolor{prussianblue}{\bfseries{C.U. BUAP}} /  \textcolor{trueblue}{\bfseries{FCFM}} \\
	\textcolor{prussianblue}{\bfseries{Edificio}} /  \textcolor{trueblue}{\bfseries{EMA4}}    & 	\textcolor{prussianblue}{\bfseries{FM9}} /  \textcolor{trueblue}{\bfseries{302}}   
\end{tabular} 
\end{mybox} \vspace{1cm}

\tableofcontents

\vspace{2cm}

Vamos a continuar con el curso, habían unas anotaciones de armónicos esféricos, entonces hoy vamos a ver otra cosa.


Tenemos el principio de Hamilton MOdificado, tenenemos Newton y que aburrido, teníamos planos inclinados en secundaria, luego en preparatoria planos inclinados con máquina de atwood. Ahora en universidad cálcula el lagrangiano, pero no sabemos cómo se calcula eso. Es imprescindible ese aspecto.

Si nosotros no empezamos a divertirnos a física, entonces no la vamos a hacer nunca. Cuando nos percatamos de hacer bien las cosas. 

Origoza tiene 58 años, no podemos pasar este día todo triste, todo aburrido. Esos días vividos tristemente. Quizá lo veamos así porque tenemos mucha edad.

Mucha de nuestra energía la estamos discipando, menos en resolver la tarea.

\fbox{Para que podamos ver cosas de física, debemos de estar atentos}. 


\section{Principio de Hamilton modificado}

Entonces lo que estamos viendo es...

Hemos partido de la segunda ley de Newton, luego de acá hemos agregado constricciones. Es algo que existe una fuerza $i$ que no podemos describir. Si la conocieramos la meteríamos así dentro de la fuerza, esto nos lleva a decir que tenemos coordenadas generalizadas que incluyen las contricciones.


\begin{equation}
\begin{split}
 m \frac{d^{r}\vec{r} }{dt^{2}} &= \vec{F} + \text{constricción}, \quad (q_{i})
\end{split}
\end{equation}


Tenemos:


\[ m \frac{d^{2} \vec{r} }{dt^{2}} = \vec{F} + \vec{F}_{c}   \]



Tenemos:


\[  \left( m \frac{ d^{2} \vec{r} }{dt^{2}} -\vec{F}  \right)  \cdot \delta{\vec{r}}  = \vec{F}_{c}\cdot \delta{\vec{r}} =0  \]



Mi partícula solo puede moverser sobre la superficie, solo tiene a la superficie como puntos accesibles. Y esa información está dentro de $\vec{F}_{c}$


\subsection{¿Qué está haciendo D'Alambert ahí?}

Eso sigue siendo cierto aún cuando la fuerza $\vec{F}$ sigue siendo cero. Mediante este procedimiento que hace. Hace magia, porque observa lo siguiente.


SI valo cero $\vec{F}$ es partícula libre. La sometemos a una esfera, resolvemos para la fuerza de constricción $\vec{F}_{c}$.

Entonces podemos construir la fuerza que nosotros queramos, Es una máquinita para crear fuerzas, por lo tanto la energía cinética, todo esto tiene que ver con nuestro curso, eso es escencial, es fundamental.


\subsection{¿Cuántas expresiones de fuerza distintas conocemos?}

En realidad son pocos, hooke, newton, fluidos, con esto anterior tenemos toda una maquinaria. 

Resulta que cuando vemos cuántica estas fuerzas pueden crearse mediante la superposición de ondas. 

\subsection{Continuando}

Tenemos el lagrangiano:

\[ L = T-V \]

Hemos dicjho que la lagrangiana depende de:


\[ L = L(q_{i},\dot{q}_{i},t)  \]

Y además hemos introducido la transformada de legendre:

y obtenemos el Hamiltoniano:


\[ H = H(q_{i}, p_{i},t)       \]

Hemos introducido tambień el potencial generalizado, Además también hemos dicho que la lagrangiana las ecuaciones de evolución:


\[ \frac{d}{dt} \left(  \frac{\partial L}{\partial \dot{q}_{i}} \right)-\frac{\partial L}{\partial q_{i}}  \]

Ecuaciones diferentciales ordinarias de segundo orden $(2n)$, pero tenemos para las coordenadas $q_{i}$. La solución involucra $2n$ constantes de integración.

Las ecuaciones de Hamilton son ecuaciones de evolución para las coordenadas $q_{i}$ y $p_{i}$:

\[ \dot{q}_{i} = \frac{\partial H }{\partial p_{i}}; \quad \dot{p}_{i} = - \frac{\partial H}{\partial q_{i}}  \]

Con estas constanes de integración, están asociadas mis contantes de movimiento, se pueden escribir en términos de las condiciones iniciales. Esas constantes de integración que surgen, en física tienen que ver con constantes de movimiento.




Si en el lagrangiano no aparece una coordenada, entoces en el lagrangiano también y esto se aplica inversamente.

Entonces todos nuestros teoremas de conservación siguen siendo válidos en esta formulación Hamiltoniana.


\subsection{Derivada total de la hamiltoniana respecto del tiempo}


Tenemos:




% Paquetes (por si acaso)
% \usepackage{amsmath,amssymb}
% \usepackage{graphicx}
% \usepackage{float}

\[
  \frac{dH}{dt} 
  = \sum_{i=1}^{\overbrace{n}^{\text{grados de libertad}}}
  \left(
    \frac{\partial H}{\partial q_{i}}\,\dot{q}_{i} 
    + 
    \frac{\partial H}{\partial p_{i}}\,\dot{p}_{i}
  \right)
  + 
  \frac{\partial H}{\partial t}
\]
Aplicando las ecuaciones de Hamilton,
\[
\dot{q}_i = \frac{\partial H}{\partial p_i},
\qquad
\dot{p}_i = -\,\frac{\partial H}{\partial q_i},
\]
los términos se cancelan y queda
\[
\boxed{\frac{dH}{dt} = \frac{\partial H}{\partial t}}
\]
Por lo tanto, si \( \frac{\partial H}{\partial t}=0 \), entonces \( \frac{dH}{dt}=0 \) y \( H \) es constante de movimiento.

\begin{figure}[H]
    \centering
    \includegraphics[width=0.7\textwidth]{1.png}
    \caption{Desarrollo esquemático en el pizarrón de la derivada temporal total de la función Hamiltoniana 
    \(H(q_i,p_i,t)\). Se parte de la expresión general 
    \(\frac{dH}{dt} = \sum_i \left( \frac{\partial H}{\partial q_i}\dot{q}_i + \frac{\partial H}{\partial p_i}\dot{p}_i \right) + \frac{\partial H}{\partial t}\), 
    y al aplicar las ecuaciones de Hamilton 
    \(\dot{q}_i = \frac{\partial H}{\partial p_i}\) y 
    \(\dot{p}_i = -\frac{\partial H}{\partial q_i}\), 
    los términos dentro de la suma se cancelan, resultando en 
    \(\frac{dH}{dt} = \frac{\partial H}{\partial t}\). 
    Si \(H\) no depende explícitamente del tiempo, se concluye que el Hamiltoniano es constante de movimiento.}
    \label{fig:derivada_total_H}
\end{figure}



Movimiento de una partícula cargada en un campo magnético uniforme. 

La lagrangiana de una partícula q n un campo magnpetico es fundamental. Esos problemas son importantes. Lagrangiano del oscilador armónico, partícula libre, problema de kpler, pero cuando tiene carga, en la tarea viene el atómo de hidrogeno. 

\end{document}